% Generated draft: requires mathematical and source review.
\documentclass[11pt]{article}
\usepackage[margin=25mm]{geometry}
\usepackage{fontspec}
\setmainfont{Latin Modern Roman}
\usepackage{amsmath,amssymb,mathtools,unicode-math}
\setmathfont{Latin Modern Math}
\usepackage{enumitem,needspace,xurl,hyperref}
\setlength{\parindent}{0pt}
\setlength{\parskip}{5pt}
\setlength{\emergencystretch}{3em}
\newcommand{\N}{\mathbb{N}}
\newcommand{\Z}{\mathbb{Z}}
\newcommand{\Q}{\mathbb{Q}}
\newcommand{\R}{\mathbb{R}}
\newcommand{\C}{\mathbb{C}}
\newcommand{\F}{\mathbb{F}}
\newcommand{\A}{\mathbb{A}}
\newcommand{\PP}{\mathbb{P}}
\newcommand{\E}{\mathbb{E}}
\newcommand{\eps}{\varepsilon}
\newcommand{\dd}{\,\mathrm{d}}
\begin{document}
\section*{UnsolvedMath ID: 30002755}
\textbf{Source problem code:} catalogue entry
\subsection{Definitions and mathematical statement}
The supplied question is: ``Does diffusion determine a domain? Equivalently, under suitable hypotheses, if two domains have the same diffusion dynamics or heat semigroup, must they be congruent?'' The hypotheses and the meaning of ``same'' are not specified, so this does not yet determine a single mathematical proposition.

A conditional formulation in a Euclidean setting is as follows; this setting is an interpretation requiring confirmation. Fix an integer \(n\geq 1\) and a class \(\mathcal D\) of nonempty open connected subsets of \(\R^n\). For each \(\Omega\in\mathcal D\), specify a function space \(X_\Omega\) and an operator \(A_\Omega\) generating the intended heat semigroup
\[
T_\Omega(t)=e^{tA_\Omega},\qquad t\geq 0.
\]
The differential operator, its domain, and its boundary conditions are part of these specifications. Two such domains are congruent if there exist \(Q\in\R^{n\times n}\) and \(b\in\R^n\) such that
\[
Q^{\mathsf T}Q=I_n,
\qquad
\Omega_2=\{Qx+b:x\in\Omega_1\}.
\]
For spaces associated with different domains, comparison of semigroups requires a specified identification. One possible comparison uses an invertible linear map \(U:X_{\Omega_1}\to X_{\Omega_2}\) satisfying
\[
U T_{\Omega_1}(t)=T_{\Omega_2}(t)U
\qquad\text{for every }t\geq 0.
\]
The admissible class of maps \(U\), including any continuity or order requirements, must be fixed. If the spaces carry an order, an order isomorphism means a linear bijection \(U\) for which both \(U\) and \(U^{-1}\) preserve nonnegativity.

With these choices fixed, the proposed implication would be
\[
\begin{gathered}
\forall\,\Omega_1,\Omega_2\in\mathcal D,\\
\bigl[\exists\text{ admissible }U:X_{\Omega_1}\to X_{\Omega_2}
\text{ with }UT_{\Omega_1}(t)=T_{\Omega_2}(t)U
\text{ for all }t\geq0\bigr]\\
\Longrightarrow\quad \Omega_1\text{ and }\Omega_2\text{ are congruent}.
\end{gathered}
\]
This is a formulation schema, not an assertion that the unspecified choices make the implication true. The source does not select the dimension, domain class, function spaces, boundary conditions, or admissible identifications. It also leaves unclear whether the intended task is to establish an implication under particular hypotheses or to characterize all hypotheses under which it holds.

The supplied dated curation describes a restricted positive result for Lipschitz domains whose heat semigroups are intertwined by an order isomorphism, and contrasts this with failure of determination by ordinary isospectrality. In a setting with discrete eigenvalue lists, isospectrality means equality of those lists, including multiplicities, for the chosen operator realizations; it is a different comparison from the specified order-isomorphism condition. The curation does not supply the full hypotheses of either cited result, and its assessment is not asserted here as an independently verified theorem.
\subsection{Short English statement}
The question asks whether a domain's diffusion determines its geometric shape up to congruence. A precise statement requires a domain class, a heat operator with boundary conditions, function spaces, and a definition of equality of diffusion across different domains. The supplied curation labels the problem partially solved and mentions a positive result using order isomorphisms and a negative comparison using isospectrality. Those claims and their exact hypotheses require human review.
\subsection{Sources}
\begin{itemize}
\item Local registry: supplied complete UnsolvedMath record, permanent ID \(30002755\), source code \texttt{catalogue entry}, title ``Determination of Domains by Diffusion,'' in the Oberwolfach Reports Open Problems collection. The record gives proposed year \(2015\) and supplies no individual proposer.
\item Supplied report attribution: \emph{Spectral Theory and Weyl Functions} (2015), Oberwolfach Reports, pp.\ 5--86, DOI \texttt{10.4171/owr/2015/1}. \url{https://doi.org/10.4171/owr/2015/1}. The background describes this report as documentation of a mathematical workshop and its open-problem discussions; no problem-specific page is supplied.
\item Supplied literature reference: Arendt, ``Does diffusion determine the body.'' \url{https://doi.org/10.1515/crll.2002.076}.
\item Supplied literature reference: Gordon, Webb, and Wolpert, ``isospectral drums.'' \url{https://doi.org/10.1090/S0273-0979-1992-00289-6}.
\item Supplied literature reference labeled ``original OWR report'': \url{https://doi.org/10.4171/owr/2015/1}.
\item Literature assessment dated exactly \texttt{2026-08-21}: the supplied curation labels the status \texttt{partially\_solved} and states that order-isomorphism intertwining gives a restricted positive determination result for Lipschitz domains, whereas ordinary isospectrality does not determine the domain. This is supplied dated curation, not independently verified here. The background explicitly calls for independent verification of time-sensitive status and cited claims.
\end{itemize}
\end{document}
